% GENERATED FILE. Edit canonical lessons, then run npm run generate:books.
% canonical-source-hash: fnv1a64:aa54c73a6a3be5b3
% canonical-lessons: TE-C435-medhavi, TE-C435-panditudu, TE-C435-tatvasastram, TE-C435-samskarta, TE-C435-prerepincu

\chapter{Thinkers and Reformers}
\label{ch:te-thinkers-and-reformers}
\hlchaptermodality{435}

\begin{chapteropening}
\textbf{By the end of this chapter:} \emph{I can talk about intellectuals, scholars, philosophy and reform.}

Complete the last lesson of chapter 435: I can talk about intellectuals, scholars, philosophy and reform.
\end{chapteropening}

\section[\texorpdfstring{\te{మేధావి} (mēdhāvi)}{మేధావి (mēdhāvi)}]{\te{మేధావి} (mēdhāvi) — an intellectual; a thinker}

\label{lesson:TE-C435-medhavi}



\begin{quote}
Before the new one: say the Telugu for an occasion, then the Telugu for to be formed.
\end{quote}

\subsection*{You'll want to know: \te{మేధావి}}
\textbf{\te{మేధావి}} — \emph{mēdhāvi} — \textquotedblleft{}an intellectual; a thinker\textquotedblright{}.

\begin{quote}
\te{ఆ} \te{సమావేశానికి} \te{చాలా} \te{మంది} \te{మేధావులు} \te{వచ్చారు}.
\end{quote}

\begin{quote}
\emph{ā samāvēśāniki cālā mandi mēdhāvulu vaccāru.}
\end{quote}

\begin{quote}
Many intellectuals came to that meeting.
\end{quote}

Plural: \te{మేధావులు}, \emph{mēdhāvulu}.

\te{మేధ}, intellect.

\begin{cousinweb}
\textbf{\te{మేధావి}} (\emph{mēdhāvi}) is borrowed from Sanskrit \emph{medhāvin}, \textquotedblleft{}intelligent; a learned man, a teacher\textquotedblright{} (Monier-Williams).
\end{cousinweb}

\begin{grammarlens}[title={What you've built}]
One more word for saying what happens if.
\end{grammarlens}

\subsection*{Your turn}
\begin{itemize}
\raggedright
  \item \emph{Say it:} \emph{mēdhāvi}
  \item \emph{Say it:} \emph{mēdhāvi}, once more
  \item \emph{Say it:} \emph{ā samāvēśāniki cālā mandi mēdhāvulu vaccāru.} — Many intellectuals came to that meeting.
  \item \emph{Recall:} say the Telugu for an occasion, then the Telugu for to be formed, then say \emph{mēdhāvi} again
\end{itemize}

\begin{tcolorbox}[breakable,colback=teal!4,colframe=teal!35!black,title={Before you move on}]
What does \te{మేధావి} mean? (\textquotedblleft{}An intellectual; a thinker\textquotedblright{}.) Say it once more.
\end{tcolorbox}
\section[\texorpdfstring{\te{పండితుడు} (paṇḍituḍu)}{పండితుడు (paṇḍituḍu)}]{\te{పండితుడు} (paṇḍituḍu) — a scholar; a learned man}

\label{lesson:TE-C435-panditudu}



\begin{quote}
Before the new one: say the Telugu for to be formed, then the Telugu for an intellectual.
\end{quote}

\subsection*{You'll want to know: \te{పండితుడు}}
\textbf{\te{పండితుడు}} — \emph{paṇḍituḍu} — \textquotedblleft{}a scholar; a learned man\textquotedblright{}.

\begin{quote}
\te{ఆ} \te{పండితుడు} \te{పాత} \te{తెలుగు} \te{బాగా} \te{చదివాడు}.
\end{quote}

\begin{quote}
\emph{ā paṇḍituḍu pāta telugu bāgā cadivāḍu.}
\end{quote}

\begin{quote}
That scholar has read a great deal of old Telugu.
\end{quote}

Plural: \te{పండితులు}, \emph{paṇḍitulu}.

A woman scholar is \te{పండితురాలు}.

\begin{cousinweb}
\textbf{\te{పండితుడు}} (\emph{paṇḍituḍu}) is borrowed from Sanskrit \emph{paṇḍita}, \textquotedblleft{}learned, clever\textquotedblright{}, and as a noun \textquotedblleft{}a scholar, a teacher\textquotedblright{} (CDIAL 7718), with the Telugu masculine ending -\te{డు}; CDIAL finds no certain etymology for the Sanskrit word.
\end{cousinweb}

\begin{grammarlens}[title={What you've built}]
One more word for saying what happens if.
\end{grammarlens}

\subsection*{Your turn}
\begin{itemize}
\raggedright
  \item \emph{Say it:} \emph{paṇḍituḍu}
  \item \emph{Say it:} \emph{paṇḍituḍu}, once more
  \item \emph{Say it:} \emph{ā paṇḍituḍu pāta telugu bāgā cadivāḍu.} — That scholar has read a great deal of old Telugu.
  \item \emph{Recall:} say the Telugu for to be formed, then the Telugu for an intellectual, then say \emph{paṇḍituḍu} again
\end{itemize}

\begin{tcolorbox}[breakable,colback=teal!4,colframe=teal!35!black,title={Before you move on}]
What does \te{పండితుడు} mean? (\textquotedblleft{}A scholar; a learned man\textquotedblright{}.) Say it once more.
\end{tcolorbox}
\section[\texorpdfstring{\te{తత్వశాస్త్రం} (tatvaśāstraṁ)}{తత్వశాస్త్రం (tatvaśāstraṁ)}]{\te{తత్వశాస్త్రం} (tatvaśāstraṁ) — philosophy}

\label{lesson:TE-C435-tatvasastram}



\begin{quote}
Before the new one: say the Telugu for an intellectual, then the Telugu for a scholar.
\end{quote}

\subsection*{You'll want to know: \te{తత్వశాస్త్రం}}
\textbf{\te{తత్వశాస్త్రం}} — \emph{tatvaśāstraṁ} — \textquotedblleft{}philosophy\textquotedblright{}.

\begin{quote}
\te{ఆమె} \te{విశ్వవిద్యాలయంలో} \te{తత్వశాస్త్రం} \te{చదువుతోంది}.
\end{quote}

\begin{quote}
\emph{āme viśvavidyālayaṁlō tatvaśāstraṁ caduvutōndi.}
\end{quote}

\begin{quote}
She is studying philosophy at university.
\end{quote}

\te{తత్వ}, the truth of things, + \te{శాస్త్రం}, a science.

\begin{cousinweb}
\textbf{\te{తత్వశాస్త్రం}} (\emph{tatvaśāstraṁ}) is built from Sanskrit \emph{tattva}, \textquotedblleft{}the real state, truth\textquotedblright{} (CDIAL 5642), and \emph{śāstra}, \textquotedblleft{}a treatise, a science\textquotedblright{} (CDIAL 12422).
\end{cousinweb}

\begin{grammarlens}[title={What you've built}]
One more word for saying what happens if.
\end{grammarlens}

\subsection*{Your turn}
\begin{itemize}
\raggedright
  \item \emph{Say it:} \emph{tatvaśāstraṁ}
  \item \emph{Say it:} \emph{tatvaśāstraṁ}, once more
  \item \emph{Say it:} \emph{āme viśvavidyālayaṁlō tatvaśāstraṁ caduvutōndi.} — She is studying philosophy at university.
  \item \emph{Recall:} say the Telugu for an intellectual, then the Telugu for a scholar, then say \emph{tatvaśāstraṁ} again
\end{itemize}

\begin{tcolorbox}[breakable,colback=teal!4,colframe=teal!35!black,title={Before you move on}]
What does \te{తత్వశాస్త్రం} mean? (\textquotedblleft{}Philosophy\textquotedblright{}.) Say it once more.
\end{tcolorbox}
\section[\texorpdfstring{\te{సంస్కర్త} (saṁskarta)}{సంస్కర్త (saṁskarta)}]{\te{సంస్కర్త} (saṁskarta) — a reformer}

\label{lesson:TE-C435-samskarta}



\begin{quote}
Before the new one: say the Telugu for a scholar, then the Telugu for philosophy.
\end{quote}

\subsection*{You'll want to know: \te{సంస్కర్త}}
\textbf{\te{సంస్కర్త}} — \emph{saṁskarta} — \textquotedblleft{}a reformer\textquotedblright{}.

\begin{quote}
\te{ఆ} \te{సంస్కర్త} \te{ఆడపిల్లల} \te{చదువు} \te{కోసం} \te{పోరాడాడు}.
\end{quote}

\begin{quote}
\emph{ā saṁskarta āḍapillala caduvu kōsaṁ pōrāḍāḍu.}
\end{quote}

\begin{quote}
That reformer fought for girls’ education.
\end{quote}

Plural: \te{సంస్కర్తలు}, \emph{saṁskartalu}.

Of the same family as \te{సంస్కరణ}, a reform.

\begin{cousinweb}
\textbf{\te{సంస్కర్త}} (\emph{saṁskarta}) is borrowed from Sanskrit \emph{saṁskartṛ}, \textquotedblleft{}one who prepares or consecrates\textquotedblright{} (Monier-Williams); Telugu uses it for one who reforms society.
\end{cousinweb}

\begin{grammarlens}[title={What you've built}]
One more word for saying what happens if.
\end{grammarlens}

\subsection*{Your turn}
\begin{itemize}
\raggedright
  \item \emph{Say it:} \emph{saṁskarta}
  \item \emph{Say it:} \emph{saṁskarta}, once more
  \item \emph{Say it:} \emph{ā saṁskarta āḍapillala caduvu kōsaṁ pōrāḍāḍu.} — That reformer fought for girls’ education.
  \item \emph{Recall:} say the Telugu for a scholar, then the Telugu for philosophy, then say \emph{saṁskarta} again
\end{itemize}

\begin{tcolorbox}[breakable,colback=teal!4,colframe=teal!35!black,title={Before you move on}]
What does \te{సంస్కర్త} mean? (\textquotedblleft{}A reformer\textquotedblright{}.) Say it once more.
\end{tcolorbox}
\section[\texorpdfstring{\te{ప్రేరేపించు} (prērēpiñcu)}{ప్రేరేపించు (prērēpiñcu)}]{\te{ప్రేరేపించు} (prērēpiñcu) — to inspire; to prompt}

\label{lesson:TE-C435-prerepincu}



\begin{quote}
Before the new one: say the Telugu for philosophy, then the Telugu for a reformer.
\end{quote}

\subsection*{You'll want to know: \te{ప్రేరేపించు}}
\textbf{\te{ప్రేరేపించు}} — \emph{prērēpiñcu} — \textquotedblleft{}to inspire; to prompt\textquotedblright{}.

\begin{quote}
\te{ఆమె} \te{మాటలు} \te{నన్ను} \te{చదువుకోవడానికి} \te{ప్రేరేపించాయి}.
\end{quote}

\begin{quote}
\emph{āme māṭalu nannu caduvukōvaḍāniki prērēpiñcāyi.}
\end{quote}

\begin{quote}
Her words inspired me to study.
\end{quote}

Past (I …): \te{ప్రేరేపించాను}, \emph{prērēpiñcānu}.

The noun \te{ప్రేరణ}, inspiration, you know.

\begin{cousinweb}
\textbf{\te{ప్రేరేపించు}} (\emph{prērēpiñcu}) is built on Sanskrit \emph{preraṇa}, \textquotedblleft{}setting in motion, urging, inciting\textquotedblright{} (Monier-Williams), with the Telugu verb ending -\te{ఇంచు}; \te{ప్రేరణ} is the same word.
\end{cousinweb}

\begin{grammarlens}[title={What you've built}]
One more word for saying what happens if.
\end{grammarlens}

\subsection*{Your turn}
\begin{itemize}
\raggedright
  \item \emph{Say it:} \emph{prērēpiñcu}
  \item \emph{Say it:} \emph{prērēpiñcu}, once more
  \item \emph{Say it:} \emph{āme māṭalu nannu caduvukōvaḍāniki prērēpiñcāyi.} — Her words inspired me to study.
  \item \emph{Recall:} say the Telugu for philosophy, then the Telugu for a reformer, then say \emph{prērēpiñcu} again
\end{itemize}

\begin{tcolorbox}[breakable,colback=teal!4,colframe=teal!35!black,title={Before you move on}]
What does \te{ప్రేరేపించు} mean? (\textquotedblleft{}To inspire; to prompt\textquotedblright{}.) Say it once more.
\end{tcolorbox}
